\documentclass[a4paper, 11pt]{article}
\usepackage[utf8]{inputenc}
\usepackage[
  left = 20mm, right = 20mm, top = 22mm, bottom = 22mm,
]{geometry}

\usepackage{amsmath,amsfonts,amssymb,bm}
\usepackage[english]{babel}
\usepackage{natbib}
\usepackage{hyperref}
\usepackage{xcolor}
\usepackage{fancyhdr}
\usepackage{booktabs}

\definecolor{harvardcrimson}{HTML}{8C1D40}

\hypersetup{
  colorlinks=true,
  linkcolor=harvardcrimson,
  citecolor=harvardcrimson,
  urlcolor=harvardcrimson,
  hypertexnames=false,
}

\pagestyle{fancy}
\fancyhf{}
\rhead{\textcolor{harvardcrimson}{\textbf{\(\Delta\Sigma^{\rm prj}\): halo-model derivation}}}
\lhead{\textcolor{harvardcrimson}{\textbf{J.\,H.\,Esteves}}}
\cfoot{\thepage}
\renewcommand{\headrulewidth}{0.4pt}
\renewcommand{\footrulewidth}{0.4pt}
\setlength{\headheight}{14pt}

\renewcommand{\baselinestretch}{1.15}
\setlength{\emergencystretch}{3em}
\sloppy

\newcommand{\lob}{\lambda^{\rm ob}}
\newcommand{\zob}{z^{\rm ob}}
\newcommand{\zcls}{z_{\rm cls}}
\newcommand{\Mcls}{M_{\rm cls}}
\newcommand{\xinl}{\xi_{\rm NL}}
\newcommand{\xilin}{\xi_{\rm lin}}
\newcommand{\xihh}{\xi_{hh}}
\newcommand{\hMpc}{h^{-1}\,\mathrm{Mpc}}
\newcommand{\hMsun}{h^{-1}\,M_{\odot}}
\newcommand{\bsel}{b_{\rm sel}}
\newcommand{\bhalo}{b_{\rm halo}}
\newcommand{\bcls}{b_{\rm cls}}
\newcommand{\Sprj}{\Sigma^{\rm prj}}
\newcommand{\Smis}{\Sigma_{\rm mis}}
\newcommand{\DSmis}{\Delta\Sigma_{\rm mis}}
\newcommand{\DSprj}{\Delta\Sigma^{\rm prj}}
\newcommand{\avg}[1]{\langle #1 \rangle}
\newcommand{\Rmis}{R_{\rm mis}}
\newcommand{\Rexcl}{R_{\rm excl}}
\newcommand{\texcl}{\theta_{\rm excl}}
\newcommand{\dO}{d\Omega}
\newcommand{\chipar}{\chi_\parallel}

\title{\textcolor{harvardcrimson}{\textbf{Halo-model derivation of \(\avg{\DSprj(R\mid\lob,\zob)}\):\\
can we avoid a second integral?}}}
\author{Johnny H.\ Esteves}
\date{\today}

\begin{document}
\maketitle
\thispagestyle{fancy}

\begin{abstract}
\noindent
We revisit the two-halo projected surface density
\(\avg{\Sprj}\) of \citet{Costanzi2026} and derive it as the
cylindrical specialisation of the Cooray--Sheth halo-halo two-halo
term (\citet{CoorayShethHaloModel}, Eqs.~86--87). We then ask a
practical question: to get \(\avg{\DSprj(R)}\), do we need a new
integral, or can we reuse the \(\avg{\Sprj(R)}\) machinery? The
answer is that the radial-average and subtraction
\(\bar\Sigma(<R)-\Sigma(R)\) is a linear functional in \(R\) only,
so it commutes with the \((\theta, z, M)\) integrals: the correct
recipe is to swap \(\Smis\!\to\!\DSmis\) inside the Eq.~13 kernel
and leave everything else alone. We also argue that, because
\(\DSmis(R'\mid\Rmis)\to 0\) for \(\Rmis\gg R'\) (the
surface-density contrast of a distant off-centre NFW neighbour
vanishes), the \(\theta\)-upper limit of the integral can be
reduced from the \(R_{\max}\!=\!30\,\hMpc\) legacy value to a
physics-driven \(\theta_{\max}\sim\text{a few}\times R'_{\max}/\chi(\zcls)\),
set by the support of \(\DSmis\) rather than by \(\xinl\) decay
or by the \texttt{twoD\_prj\_NFW} truncation.
\end{abstract}

\section{Setup: Cooray--Sheth two-halo term in cylindrical coordinates}
\label{sec:setup}

\paragraph{Geometry and notation.}
Fix the cluster at the origin, at comoving distance
\(\chi(\zcls)\); its mass is \(\Mcls\) and its halo bias
\(\bcls\equiv b(\Mcls,\zcls)\). A neighbouring halo (the
two-halo contributor, ``2h'') sits at sky angle \(\theta\) from
the cluster and redshift \(z\), with mass \(M\) and halo bias
\(b(M,z)\). We adopt cylindrical coordinates
\((R,\,\chipar)\) around the cluster's line of sight, with
\begin{equation}
R \equiv \theta\,\chi(z) \;\approx\; \theta\,\chi(\zcls) ,
\qquad
\chipar \equiv \chi(z) - \chi(\zcls) ,
\qquad
r \equiv \sqrt{R^2+\chipar^2} ,
\label{eq:cyl_coords}
\end{equation}
where the small-angle approximation \(R\simeq\theta\,\chi(\zcls)\)
holds across the photo-\(z\) window (\(|\chipar|\ll\chi(\zcls)\))
and \(r\) is the 3-D comoving cluster--neighbour separation. The
exact form used numerically
(\texttt{sigma\_prj.py}, \S3 of the refactor note) is the chord
\begin{equation}
r^2 \;=\; \chi(z)^2 + \chi(\zcls)^2 - 2\,\chi(z)\,\chi(\zcls)\cos\theta ,
\label{eq:r_exact}
\end{equation}
which reduces to~\eqref{eq:cyl_coords} at \(\theta\ll 1\).

\paragraph{Starting point: Cooray--Sheth Eq.~86.}
The 3-D matter two-halo correlator is
\citep[Eq.~86]{CoorayShethHaloModel}
\begin{multline}
\xi_{2h}(\mathbf{x}-\mathbf{x}')
= \int\!d\Mcls\,\tfrac{\Mcls n(\Mcls)}{\bar\rho}
  \int\!dM\,\tfrac{M\,n(M)}{\bar\rho} \\
  \times \int\!d^3x_c\,u(\mathbf{x}-\mathbf{x}_c\mid\Mcls)
  \int\!d^3x\,u(\mathbf{x}'-\mathbf{x}\mid M)\,
  \xihh(\mathbf{x}_c-\mathbf{x}\mid\Mcls,M),
\label{eq:CS86}
\end{multline}
with \(\mathbf{x}_c\) the cluster centre, \(\mathbf{x}\) the 2h
halo centre, and \(u(\cdot\mid M)\) the unit-normalised density
profile. For our observable we fix the cluster centre at the
origin, \(\mathbf{x}_c\equiv\mathbf{0}\), and drop its own profile
(the 1-halo term is treated separately). What remains is the
surface density around the cluster contributed by projected
neighbouring halos:
\begin{equation}
\Sprj(R\mid\Mcls,\zcls)
= \int d^2R_\perp \int d\chipar \int dM\,n(M,z)\,
  \Sigma_{\rm NFW}\!\bigl(|\mathbf{R}-\mathbf{R}_\perp|\,\big|\,M,z\bigr)\,
  \xihh\!\bigl(r\,\big|\,\Mcls,M\bigr),
\label{eq:cyl_raw}
\end{equation}
where \(\mathbf{R}_\perp\) is the transverse offset of the
neighbour from the cluster line of sight (\(|\mathbf{R}_\perp|=R\)),
\(\chipar\) is its LoS offset, and
\(\Sigma_{\rm NFW}(R'\mid M,z)\) is the centred NFW surface
density of the neighbour evaluated at 2-D radius \(R'\) from its
own centre.

\paragraph{Change of coordinates to \((\theta, z)\).}
Swap the neighbour's \((\mathbf{R}_\perp,\chipar)\) for the
observable pair \((\theta, z)\). The volume element becomes
\begin{equation}
d^2R_\perp\,d\chipar
\;=\; 2\pi\sin\theta\,d\theta \,\cdot\, \frac{dV}{dz\,\dO}(z)\,dz .
\label{eq:volume_element}
\end{equation}
Azimuth-averaging \(\Sigma_{\rm NFW}(|\mathbf{R}-\mathbf{R}_\perp|)\)
around the neighbour's off-centre offset \(\Rmis\equiv R\)
produces the \emph{miscentered} NFW surface density,
\begin{equation}
\Smis(R'\mid M,z,\Rmis)
= \int_0^{2\pi}\!d\varphi\;
  \Sigma_{\rm NFW}\!\Bigl(\sqrt{R'^{\,2}+\Rmis^2-2R'\Rmis\cos\varphi}
  \,\Big|\,M,z\Bigr),
\label{eq:Smis_def}
\end{equation}
where \(R'\) is the radius at which the cluster observer measures
\(\Sprj\). For the cluster--neighbour geometry
\(\Rmis = \theta\,\chi(\zcls)\) (equation~\ref{eq:cyl_coords}).
Equation~\eqref{eq:cyl_raw} becomes
\begin{multline}
\Sprj(R'\mid\Mcls,\zcls) =
\int d\theta\sin\theta\int dz\,\frac{dV}{dz\,\dO}(z)
\int dM\,n(M,z) \\
\times\; \xihh\!\bigl(r\mid\Mcls,M\bigr)\,
\Smis\!\bigl(R'\mid M,z,\Rmis=\theta\,\chi(\zcls)\bigr),
\label{eq:two_halo_cyl}
\end{multline}
with \(r\) given by~\eqref{eq:r_exact}.

\paragraph{Cooray--Sheth Eq.~87 step.}
Deterministic linear halo bias replaces the halo-halo correlator
by matter clustering times two bias factors,
\(\xihh(r\mid\Mcls,M)\approx\bcls\,b(M,z)\,\xilin(r)\).
Factoring \(\bcls\) out:
\begin{multline}
\Sprj(R'\mid\Mcls,\zcls) \approx \bcls
\int d\theta\sin\theta\int dz\,\frac{dV}{dz\,\dO}
\int dM\,n(M,z)\,b(M,z)\\
\times\; \xilin(r)\,\Smis\!\bigl(R'\mid M,z,\theta\,\chi(\zcls)\bigr).
\label{eq:CS87_cyl}
\end{multline}

\paragraph{From mass-selected to richness-selected.}
\citet{Costanzi2026} upgrade \eqref{eq:CS87_cyl} in three steps to
describe the observable around an optically selected cluster of
observed richness \(\lob\):
(i) \(\xilin\to\xinl\) (halofit), because the 1h--2h transition at
\(\sim\Rexcl\) sits in the non-linear regime;
(ii) \(\xinl\to 0\) for \(\theta<\texcl(z)\) (LoS-slab exclusion);
(iii) the cluster's bias \(\bcls\) is replaced by the
scale-dependent \emph{selection-affected} bias
\(\bsel(\theta;\lob,\zob)\), which absorbs red-galaxy selection
effects. Adding the uncorrelated cosmological mean as a \(+1\)
inside the bracket yields the master equation
(\citealt{Costanzi2026}, Eq.~13):
\begin{multline}
\avg{\Sprj(R'\mid\lob,\zob)} =
\int d\theta\sin\theta\int dz\,\frac{dV}{dz\,\dO}\int dM\,n(M,z)\\
\times\;\bigl[\,1 + b(M,z)\,\bsel(\theta;\lob,\zob)\,\xinl(r,\zob)\,\bigr]\,
\Smis\!\bigl(R'\mid M,z,\Rmis=\theta\,\chi(\zcls)\bigr).
\label{eq:Sprj}
\end{multline}
The \(1\) piece integrates to a spatially near-uniform
\(\Sigma_{\rm rnd}(R')\) (the cosmological mean projected surface
density in the photo-\(z\) window); the \(b\,\bsel\,\xinl\) piece
is the correlation-excess two-halo contribution
\(\Sigma_{\rm cl+LSS}(R')\).

\section{\texorpdfstring{The target observable: \(\avg{\DSprj(R')}\)}{The target observable: DeltaSigma prj}}
\label{sec:dsprj}

The lensing observable is the \emph{excess} surface density,
\begin{equation}
\avg{\DSprj(R')} \equiv \bar\Sigma^{\rm prj}(<R') - \avg{\Sprj(R')},
\qquad
\bar\Sigma^{\rm prj}(<R') \equiv \frac{2}{R'^{\,2}}\int_0^{R'}\!\!
s\,\avg{\Sprj(s)}\,ds .
\label{eq:DSprj_def}
\end{equation}

\paragraph{Linearity in \(R'\).}
Both \(\bar\Sigma(<R')\) and the subtraction act only on the
radial argument \(R'\). The outer integrals in~\eqref{eq:Sprj}
(over \(\theta, z, M\)) are \(R'\)-independent, so they commute
with \(\Delta[\,\cdot\,]\):
\begin{multline}
\avg{\DSprj(R'\mid\lob,\zob)} =
\int d\theta\sin\theta\int dz\,\frac{dV}{dz\,\dO}\int dM\,n(M,z)\\
\times\;\bigl[\,1 + b(M,z)\,\bsel(\theta;\lob,\zob)\,\xinl(r,\zob)\,\bigr]\,
\DSmis\!\bigl(R'\mid M,z,\Rmis=\theta\,\chi(\zcls)\bigr),
\label{eq:DSprj}
\end{multline}
where the only change from~\eqref{eq:Sprj} is
\begin{equation}
\Smis\!\bigl(R'\mid M,z,\Rmis\bigr)
\;\longrightarrow\;
\DSmis\!\bigl(R'\mid M,z,\Rmis\bigr)
\equiv \bar\Sigma_{\rm mis}(<R'\mid M,z,\Rmis) - \Smis(R'\mid M,z,\Rmis).
\label{eq:DSmis_def}
\end{equation}
The machinery that computes the \((\theta, z, M)\) integral (the
split-at-breakpoints \(\theta\)-grid with per-\(R'\) nodes, the
photo-\(z\) \(z\)-slab, the exact \(r\) form, the exclusion cut)
is preserved --- \(\DSmis\) just replaces \(\Smis\) as a lookup
against the offset-NFW tables.

\paragraph{Does the \(\Sigma_{\rm rnd}\) piece drop out?}
In the idealised sky-limit \(\theta_{\max}\!\to\!\infty\) with
infinitely narrow 1-halo profiles, the \(1\) piece gives a
spatially uniform \(\Sigma_{\rm rnd}\), for which
\(\bar\Sigma(<R')\equiv\Sigma_{\rm rnd}\) and thus
\(\Delta\Sigma_{\rm rnd}(R')=0\). This is the classical
\citet{Sheldon2009}--\citet{Zu2014}--\citet{Melchior2017}
statement that the mean photo-\(z\)-window surface density cancels
from the measured excess. Numerically the package truncates
\(\theta\) at \(\theta_{\max}=R_{\max}/\chi(\zcls)\) with
\(R_{\max}=30\,\hMpc\) (\texttt{config.R\_MAX\_CMPCH}), so
\(\Delta\Sigma_{\rm rnd}\) is not exactly zero but a boundary term
that scales with \(R_{\max}\). We therefore recommend, by analogy
with \(\avg{\Sprj}\), returning only the cl+LSS piece by default:
\begin{multline}
\avg{\DSprj_{\rm cl+LSS}(R')}
= \int d\theta\sin\theta\int dz\,\frac{dV}{dz\,\dO}\int dM\,n(M,z)\\
\times\; b(M,z)\,\bsel(\theta)\,\xinl(r,\zob)\,
\DSmis\!\bigl(R'\mid M,z,\theta\,\chi(\zcls)\bigr).
\label{eq:DSprj_cl}
\end{multline}

\section{\texorpdfstring{Choosing \(\theta_{\max}\) for \(\avg{\DSprj}\)}{Choosing theta_max for DeltaSigma prj}}
\label{sec:theta_max}

The \(\theta\)-upper limit of \eqref{eq:Sprj}/\eqref{eq:DSprj} is
an artefact of numerical truncation; the physics runs to
\(\theta\!\to\!\pi\). For \(\avg{\Sprj}\) the legacy choice is
\(R_{\max}=30\,\hMpc\), so that
\(\theta_{\max}=R_{\max}/\chi(\zcls)\), matching the hard
truncation in Matteo's \texttt{twoD\_prj\_NFW} notebook
(\texttt{config.R\_MAX\_CMPCH}, \S5.1 of the refactor note). At
that scale the \texttt{cl+LSS} piece is converged to \(<\!1\%\)
because \(\xinl(r)\) has decayed, but the \texttt{rnd} piece
grows by \(\sim 15\%\) going \(30\!\to\!60\,\hMpc\): the flat
large-\(\theta\) plateau of the \texttt{rnd} integrand
never dies \citep[see][\S8]{SigmaPrjRefactor}. That is why the
package returns \texttt{cl+LSS} by default for \(\avg{\Sprj}\).

\paragraph{What changes for \(\avg{\DSprj}\).}
Two facts together imply that the \(\theta_{\max}\) requirement
for \(\avg{\DSprj}\) is strictly \emph{weaker} than for
\(\avg{\Sprj}\):
\begin{enumerate}\setlength\itemsep{2pt}
  \item \textbf{The \texttt{rnd} piece is dropped by design.}
        Following~\eqref{eq:DSprj_cl}, \(\avg{\DSprj}\) in this
        package is the \texttt{cl+LSS}-only integral --- the
        long \(\theta\)-plateau of the \texttt{rnd} piece simply
        does not appear in the observable.
  \item \textbf{The \texttt{cl+LSS} \(\theta\)-integrand of
        \(\avg{\DSprj}\) is itself bounded in \(\theta\).}
        \(\DSmis(R'\mid\Rmis)\) has compact support around
        \(\Rmis\sim R'\) and decays in both limits
        (\(\Rmis\ll R'\) gives
        \(\DSmis\!\to\!2\pi\Delta\Sigma_{\rm NFW}(R')\), falling
        in \(R'\); \(\Rmis\gg R'\) gives
        \(\Smis\!\to\!2\pi\Sigma_{\rm NFW}(\Rmis)\),
        \(R'\)-independent, so
        \(\DSmis=\bar\Sigma_{\rm mis}-\Smis\!\to\!0\)). In contrast, the
        \(\avg{\Sprj}\) \texttt{cl+LSS} integrand inherits the
        slow plateau of \(\Smis(R'\mid\Rmis)\) at large
        \(\Rmis\), which in numerical plots shows up as a near-constant
        \(N_{\rm cl+LSS}(\theta)\) beyond \(\theta_\lambda\).
\end{enumerate}
In other words, the two mechanisms that forced
\(\theta_{\max}\!=\!30\,\hMpc/\chi(\zcls)\) for \(\avg{\Sprj}\)
--- the \texttt{rnd} plateau \emph{and} the flat
\(\Smis\)-kernel at large \(\Rmis\) --- both go away when we
switch to \(\DSprj_{\rm cl+LSS}\). The remaining constraint on
\(\theta_{\max}\) is only that we resolve the support of
\(\DSmis(R'\mid\theta\,\chi(\zcls))\), which is set by the
largest measurement radius \(R'_{\max}\).

\paragraph{Recipe.}
Take \(\theta_{\max}\) as the first of
\(\{\theta_{R'_{\max}}\) times a safety factor,
\(\theta_\lambda\) times a safety factor\(\}\) that is large
enough to include the tail of the ring peak at
\(\Rmis=R'_{\max}\). Empirically,
\begin{equation}
\boxed{\;
\theta_{\max}^{\Delta\Sigma}
\;=\;
\frac{C\,R'_{\max}}{\chi(\zcls)} ,
\qquad C\simeq 3 ,\;
R'_{\max}\equiv\max_i R'_i .
\;}
\label{eq:theta_max_DS}
\end{equation}
For a DES-Y3-style setup with
\(R'_{\max}\!=\!10\,\hMpc\) and \(\zcls\!=\!0.5\),
\(\theta_{\max}^{\Delta\Sigma}\!\approx\!30/\chi(\zcls)\), so
numerically it happens to coincide with the \(\avg{\Sprj}\)
default; for \(R'_{\max}\!=\!3\,\hMpc\) it drops to
\(\sim 9/\chi(\zcls)\), a \(\sim 3\times\) cheaper
\(\theta\)-integral. Do \emph{not} push \(C\) below \(\sim 2\):
the ring peak of \(\DSmis\) at \(\Rmis=R'_{\max}\) extends
modestly beyond \(\theta_{R'_{\max}}\) because the NFW neighbour
has a finite \(R_s\sim 0.3\,\hMpc\).

\paragraph{\(\xinl\)-decay is not the binding constraint.}
One might worry that \(\xinl(r)\) support extends to tens of
\(\hMpc\) and therefore forces \(\theta_{\max}\) large even for
\(\avg{\DSprj}\). It does not, for the same reason the
\texttt{rnd} plateau drops out: at large \(\theta\) the
\(\DSmis\)-kernel is what zeros the integrand, not \(\xinl\).
\(\xinl\) decay would bind \(\theta_{\max}\) only for an
observable where \(\DSmis\) is replaced by an
\(R'\)-independent kernel (e.g.\ the \(\avg{\Sprj}\) \texttt{rnd}
piece). Here, taking \(\theta_{\max}\) out past
\(\theta_{R'_{\max}}\!\times\!(\text{few})\) is wasted work.

\paragraph{Regression check.}
An analogue of \texttt{validations/sigma\_prj\_theta\_max.py}
for the \(\DSprj_{\rm cl+LSS}\) observable should show:
\begin{itemize}\setlength\itemsep{2pt}
  \item moving \(\theta_{\max}\!:\!R_{\max}=30\!\to\!60\,\hMpc\)
        leaves \(\DSprj_{\rm cl+LSS}(R')\) unchanged to
        \(\ll 1\%\) on every \(R'\!\le\!R'_{\max}\),
  \item reducing \(\theta_{\max}\!:\!R_{\max}=30\!\to\!3\,R'_{\max}\)
        (\(\approx 9\,\hMpc\) for \(R'_{\max}\!=\!3\))
        also leaves \(\DSprj_{\rm cl+LSS}\) unchanged to
        \(\ll 1\%\),
  \item reducing further (e.g.\ \(R_{\max}\!=\!R'_{\max}\))
        starts to clip the ring-peak tail at \(\Rmis=R'_{\max}\)
        and produces a \(\gtrsim 1\%\) bias at
        \(R'\!\sim\!R'_{\max}\).
\end{itemize}
The \(C\!=\!3\) floor in~\eqref{eq:theta_max_DS} is set to leave
a healthy margin before that last regime.

\section{What changes in the code}

In \texttt{sigma\_prj.py}, the only functional change needed to
compute \(\avg{\DSprj(R')}\) is to replace the kernel lookup
\begin{center}
\texttt{NFWMiscentered.sigma\_grid(R, Rtheta, M, z)}
\(\longrightarrow\)
\texttt{NFWMiscentered.delta\_sigma\_grid(R, Rtheta, M, z)}
\end{center}
where \texttt{delta\_sigma\_grid} returns
\(\bar\Sigma_{\rm mis}(<R'\mid M,z,\Rmis)-\Smis(R'\mid M,z,\Rmis)\) at the
same offset-NFW lookup resolution and with the same factor-of-2
convention as \texttt{sigma\_grid}. The second change is to make
the \(\theta\)-upper limit adaptive to the requested \(R'\)-grid
rather than fixed at \texttt{R\_MAX\_CMPCH}: set
\(R_{\max}^{\Delta\Sigma} = C\cdot R'_{\max}\) with \(C\!=\!3\)
(eq.~\ref{eq:theta_max_DS}), plumbed into the \(\theta\)-grid
constructor in place of the hard-coded \(30\,\hMpc\). Everything
else --- the \(\theta\)-outer loop, the split-at-breakpoints
\(\theta\)-grid with per-\(R'\) nodes, the
ring\,\(\cup\)\,outer-fg\,\(\cup\)\,outer-bg \(z\)-grid, \(\bsel\)
evaluation --- is preserved, and the \texttt{cl+LSS}/\texttt{rnd}
decomposition becomes moot because the default return is
\texttt{cl+LSS} only (eq.~\ref{eq:DSprj_cl}).

\begin{thebibliography}{9}
\bibitem[Cooray \& Sheth(2002)]{CoorayShethHaloModel}
Cooray, A., \& Sheth, R. 2002, Phys.\ Rep., 372, 1.
Eqs.~86--87 of the halo-model review.

\bibitem[Costanzi et al.(2026)]{Costanzi2026}
Costanzi, M., Wu, H.-Y., Esteves, J.~H., Grandis, S., To, C., \&
Aguena, M., 2026, Phys.~Rev.~D (accepted), arXiv:2604.05833.
\(\bsel\), \(\avg{\Sprj}\) (Eq.~13).

\bibitem[Esteves(2026)]{SigmaPrjRefactor}
Esteves, J.~H.\ 2026, \texttt{richness\_selection} numerical
recipe, \texttt{docs/sigma\_prj\_refactor.md}. \(\avg{\Sprj}\)
\(\theta\)-grid, \(R_{\max}\) sensitivity, \texttt{rnd}/\texttt{cl+LSS}
decomposition.

\bibitem[Sheldon et al.(2009)]{Sheldon2009}
Sheldon, E.~S., et al.\ 2009, ApJ, 703, 2217. Random-point
subtraction convention for cluster lensing.

\bibitem[Zu et al.(2014)]{Zu2014}
Zu, Y., Weinberg, D.~H., Rozo, E., Sheldon, E.~S., Tinker, J.~L.,
\& Becker, M.~R.\ 2014, MNRAS, 439, 1628.

\bibitem[Melchior et al.(2017)]{Melchior2017}
Melchior, P., et al.\ 2017, MNRAS, 469, 4899. DES Y1 cluster
lensing; uses random-point subtraction to remove \(\Sigma_{\rm rnd}\).
\end{thebibliography}

\end{document}
